% =====================================================================
% MIKU SCI 小论文 — 中文稿
% 自包含可编译骨架，使用 latexmk -xelatex main.tex 编译
% =====================================================================
\documentclass[a4paper,zihao=5]{ctexart}

\usepackage{amsmath,amssymb,amsthm}
\usepackage{booktabs}
\usepackage{graphicx}
\usepackage{enumitem}
\usepackage{geometry}
\usepackage[ruled,vlined,linesnumbered]{algorithm2e}
\DontPrintSemicolon
\SetAlgoCaptionSeparator{\quad}
\SetAlgorithmName{算法}{算法}{算法目录}
\usepackage{tikz}
\usetikzlibrary{shapes.geometric, arrows.meta, positioning, fit, backgrounds, calc, patterns, decorations.pathreplacing}
\usepackage{caption}
\captionsetup{font=small,skip=4pt,labelsep=space}
\usepackage{subcaption}
\usepackage{float}
\usepackage{adjustbox}
\usepackage{pgfplots}
\pgfplotsset{compat=1.18}
\usepgfplotslibrary{fillbetween}
\pgfplotsset{table/search path={../图片, ../图片/data}}
\providecommand{\datapath}{../图片/data/01_crossing_ped/}
% 概念图内部引用的示例参数宏，给出缺省值以保证自包含编译
\providecommand{\PedExampleS}{10.0}
\providecommand{\PedExampleTau}{1.25}
\providecommand{\PedExampleVel}{8.0}
\providecommand{\PedExampleLatVel}{1.7}
\providecommand{\PedExampleDelta}{1.2}
\providecommand{\PedExamplePredL}{2.125}
\providecommand{\PedExampleLeftBoundary}{1.875}
\providecommand{\PedExampleRightBoundary}{2.375}
\providecommand{\PedExampleGapZero}{2.55}
\providecommand{\PedExampleGapOne}{-1.70}
\providecommand{\PedExampleDeltaT}{0.1}
\providecommand{\PedExampleSk}{10.0}
\providecommand{\PedExampleTauK}{1.25}
\providecommand{\PedExampleBoundaryL}{0.0}
\providecommand{\PedExampleRoadLeft}{-1.875}
\providecommand{\PedExampleRoadRight}{1.875}
\providecommand{\PedExampleRoadWidth}{3.75}
\providecommand{\PedExamplePedWidth}{0.5}
\providecommand{\PedExamplePedHalfWidth}{0.25}
\providecommand{\PedExampleEgoWidth}{2.11}
\providecommand{\PedExampleMidTau}{0.625}
\providecommand{\PedExampleMidS}{5.0}
\providecommand{\PedExampleMidLatShift}{1.06}
\providecommand{\PedExampleFinalL}{2.125}
\providecommand{\PedExampleBoundaryLow}{1.875}
\providecommand{\PedExampleBoundaryHigh}{2.375}
\providecommand{\PedExampleBaselineGap}{0.425}
\providecommand{\PedExampleLastRestrictedIndex}{12}
\providecommand{\PedExampleLastRestrictedTime}{1.2}
\providecommand{\PedExampleFirstFreeIndex}{13}
\providecommand{\PedExampleFirstFreeTime}{1.3}
\providecommand{\PedExampleOvershoot}{0.25}

% ---- 图片目录与统一配色（复用毕设 图片/ 下自包含 TikZ 概念图）----
\graphicspath{{../图片/}}
\makeatletter
\def\input@figpath#1{\input{../图片/#1}}
\makeatother
\definecolor{deepblue}{RGB}{81, 132, 178}
\definecolor{lightblue}{RGB}{170, 212, 248}
\definecolor{skyblue}{RGB}{154, 220, 255}
\definecolor{inputyellow}{RGB}{255, 248, 154}
\definecolor{outputpink}{RGB}{255, 178, 166}
\definecolor{improvedpink}{RGB}{255, 138, 174}
\definecolor{lightpink}{RGB}{241, 167, 181}
\definecolor{deeppink}{RGB}{213, 82, 118}
\definecolor{bglight}{RGB}{242, 245, 250}
\definecolor{lightgreen}{RGB}{180, 230, 180}
\definecolor{lightorange}{RGB}{255, 210, 140}
\definecolor{dangerred}{RGB}{230, 100, 100}
\definecolor{vibrantgreen}{RGB}{46, 175, 80}
\definecolor{vibrantorange}{RGB}{240, 130, 30}
\definecolor{vibrantred}{RGB}{220, 50, 50}
\tikzset{
  mod/.style            = {rectangle, rounded corners=2pt, draw=deepblue, fill=lightblue, minimum height=0.9cm, align=center, font=\small},
  modimp/.style         = {rectangle, rounded corners=2pt, draw=deeppink, fill=improvedpink, minimum height=0.9cm, align=center, font=\small},
  modin/.style          = {rectangle, draw=deepblue, fill=inputyellow, minimum height=0.9cm, align=center, font=\small},
  modout/.style         = {rectangle, draw=deeppink, fill=outputpink, minimum height=0.9cm, align=center, font=\small},
  axx/.style            = {->, thick, black!80},
  flow/.style           = {-Stealth, thick, deepblue},
  flowimp/.style        = {-Stealth, thick, deeppink},
  obsstatic/.style      = {fill=dangerred!70, draw=dangerred, thick},
  obsdyn/.style         = {fill=lightorange, draw=dangerred, thick},
  obsbuf/.style         = {draw=dangerred!70, dashed, thick},
  bandok/.style         = {fill=lightgreen, opacity=0.55},
  bandbad/.style        = {fill=dangerred!30, opacity=0.55},
  egopt/.style          = {fill=deepblue, draw=deepblue},
  reflinestyle/.style   = {deepblue, very thick},
}
\usepackage[backend=biber,style=gb7714-2015,gbpub=false,
            uniquename=init,uniquelist=true,
            maxcitenames=2,mincitenames=1,
            sorting=none,doi=false,url=false,
            natbib=true]{biblatex}

\geometry{margin=2.2cm}
\ctexset{
  section/beforeskip   = 2.2ex plus .5ex minus .2ex,
  section/afterskip    = 1.6ex plus .3ex,
  subsection/beforeskip= 1.8ex plus .4ex minus .2ex,
  subsection/afterskip = 1.2ex plus .3ex,
}

% ---- 定理类环境（中文，按节编号）----
\theoremstyle{definition}
\newtheorem{proposition}{命题}[section]
\newtheorem{lemma}{引理}[section]
\newtheorem{theorem}{定理}[section]
\renewcommand{\proofname}{证明}

% ---- 行距与列表 ----
\linespread{1.02}
\setlist{nosep}
\setlength{\textfloatsep}{8pt plus 2pt minus 2pt}
\setlength{\intextsep}{8pt plus 2pt minus 2pt}
\setlength{\abovedisplayskip}{5pt plus 1pt minus 1pt}
\setlength{\belowdisplayskip}{5pt plus 1pt minus 1pt}
\emergencystretch=3em  % 容纳无法断行的长等宽标识符，避免 overfull

% ---- 参考文献 ----
\addbibresource{references.bib}
\renewcommand*{\bibfont}{\footnotesize}
\setlength{\bibitemsep}{1.5pt}

\title{MIKU 面向路径--速度解耦规划器的\\多障碍物时空同伦约束统一方法}
\author{廖嘉旺，孙成娇\thanks{通讯作者。}\\[2pt]\normalsize 湖北文理学院 计算机工程学院}
\date{}

\begin{document}
\maketitle

\begin{abstract}
\small
\noindent
路径--速度解耦规划器适合在线规划，但时间盲的路径边界会叠加不同时刻的动态占据，逐障碍物避让也可能因局部选向冲突而误判阻塞。本文提出多障碍物时空同伦约束统一方法 MIKU，通过交互感知安全裕度、到达时间驱动投影和组级空间同伦构造时变通行带，并将通行带成立所需的时序条件写入原速度 QP。连续分割将 $k$ 个障碍物的绕行组合化为 $k+1$ 个候选间隙，使单截面搜索由 $O(2^k)$ 降为 $O(k\log k)$，而多障碍物时间同伦则从动态占用补集中选择可达安全窗。在 4,000 组随机区间上，最大间隙扫描与组合枚举最优解逐例得到相同的单截面中心通行带宽度。在覆盖 7 类交通构型的 3,500 个配对场景中，MIKU 将无碰撞终点到达率由 62.7\% 提高至 77.5\%，配对差值的 95\% 置信区间为 $[13.5,16.0]$ 个百分点，碰撞率由 2.2\% 降至 1.1\%，同时使总体平均进度率提高 8.9 个百分点，规划耗时中位数约为 10\,ms。

\vspace{0.5em}
\noindent\textbf{关键词}\quad 自动驾驶，路径--速度解耦，时变通行带，多障碍物决策，组合优化
\end{abstract}
% ====== 1 引言 ======
\section{引言}
\label{sec:intro}

轨迹规划需要在动态环境中同时处理空间可通行性和到达时序。路径--速度解耦（Path-Velocity Decomposition, PVD）把两者分成 Frenet 空间的横向路径与固定路径上的纵向速度，因而适合在线规划\citep{kant1986pvd,werling2010,zhou2021dliaps}，也被 Apollo 等工程规划器采用\citep{fan2018baidu,apollo2024,zhou2021pathqp,wang2022_jfr}。但在路径边界构造中，动态障碍物常被简化为当前占据或预测时域的并集，多个障碍物又按局部规则依次选择绕行方向，前者会叠加本不同时发生的占据，后者会遗漏组内其他可行选向，两种机制都可能使中心可行区间过度收缩\citep{meng2019_access,gonzalez2016,paden2016}。图\ref{fig:q34_teaser}给出了逐个不可回退的选向压缩上下界，而组级连续分割保留有效间隙的情形。

联合时空搜索可以同时处理绕行与让行，但通常要更换或扩展原解耦求解链。本文在不对路径 QP 与速度 QP 的目标函数和运动学等式改造的前提下，通过重构其环境约束接口，使路径边界按照预计到达时刻和组级绕行组合生成，并将相应时序条件写入 ST 位置边界。

本文的主要工作包括
\begin{enumerate}[leftmargin=1.6em,itemsep=0.25em]
    \item 将到达时间驱动的动态 SL 投影与差异化安全裕度写成逐截面边界生成算子，使路径阶段能够区分“当前被占据”与“自车到达时仍被占据”。
    \item 用连续分割引理把组内 $2^k$ 种绕行组合压成可证明的 $k+1$ 个候选间隙，并以 $O(k\log k)$ 最大间隙扫描求得最宽中心通行带，再由保留宽度最大的 $K$ 个候选带的空间同伦保持分量之间的纵向连续。
    \item 将空间同伦所依赖的时序前提经安全窗和时间同伦图编译为标准 ST 盒约束，使路径阶段的离散决策可以直接进入原速度 QP。
\end{enumerate}

第\ref{sec:related}节归纳相关规划体系，并形式化解耦求解接口与两类可行域收缩，第\ref{sec:method}节给出 MIKU 的约束构造与复杂度，第\ref{sec:experiment}节报告统计对照、机制案例和规划链接入，第\ref{sec:conclusion}节总结。

\begin{figure}[htbp]
\centering
\begin{subfigure}[b]{0.48\textwidth}
  \centering
  \adjustbox{max width=\linewidth}{% 贪心独立决策 → BLOCKED（子图 a）
\begin{tikzpicture}[scale=0.85, transform shape]
  % 道路
  \fill[bglight] (-0.5,-2.5) rectangle (8.5,2.5);
  \draw[gray, thick] (-0.5, 2.5) -- (8.5, 2.5);
  \draw[gray, thick] (-0.5,-2.5) -- (8.5,-2.5);
  \draw[dashed, gray!60] (-0.5, 0) -- (8.5, 0);
  \node[gray, font=\scriptsize, anchor=east] at (-0.6, 0) {$l{=}0$};

  % O_a（上方靠边）y=[0.5, 1.0]
  \fill[dangerred!70, draw=dangerred, thick] (3.2, 0.5) rectangle (4.8, 1.0);
  \node[font=\scriptsize, white] at (4.0, 0.75) {$O_a$};

  % O_b（下方靠边）y=[-1.0, -0.5]
  \fill[dangerred!70, draw=dangerred, thick] (3.2,-1.0) rectangle (4.8,-0.5);
  \node[font=\scriptsize, white] at (4.0,-0.75) {$O_b$};

  % 膨胀虚线 δ=0.4
  \draw[dangerred!60, dashed, thick] (2.8, 0.1) rectangle (5.2, 1.4);
  \draw[dangerred!60, dashed, thick] (2.8,-1.4) rectangle (5.2,-0.1);

  % 决策标注
  \node[deepblue, font=\scriptsize\bfseries] at (5.5, 0.75) {R};
  \node[deeppink, font=\scriptsize\bfseries] at (2.5,-0.75) {L};

  % 独立决策：O_a 右绕 → 上界压到 0.1-0.1=0.0
  %           O_b 左绕 → 下界抬到 -0.1+0.1=0.0
  % l+ ≈ l- → BLOCKED

  % 上界 l+(s) 凹字形（offset 区分）
  \draw[deeppink, very thick]
    (-0.5, 2.5) -- (2.6, 2.5) -- (2.6, 0.0) -- (5.4, 0.0) -- (5.4, 2.5) -- (8.5, 2.5);
  \node[deeppink, font=\scriptsize] at (7.5, 2.1) {$l^+(s)$};

  % 下界 l-(s) 凸字形
  \draw[deepblue, very thick]
    (-0.5,-2.5) -- (2.6,-2.5) -- (2.6, 0.0) -- (5.4, 0.0) -- (5.4,-2.5) -- (8.5,-2.5);
  \node[deepblue, font=\scriptsize] at (7.5,-2.1) {$l^-(s)$};

  % BLOCKED 标注
  \node[dangerred, font=\scriptsize\bfseries] at (4.0, 1.8) {$l^+ \leq l^-$ $\Rightarrow$ BLOCKED};

  % 自车
  \fill[deepblue, opacity=0.4] (0.3,-0.3) rectangle (1.8,0.3);
  \draw[deepblue, thick, rounded corners=1pt] (0.3,-0.3) rectangle (1.8,0.3);
  \node[deepblue, font=\tiny\bfseries] at (1.05,0) {自车};
\end{tikzpicture}
}
  \caption{逐障碍物选向产生虚假阻塞}
\end{subfigure}\hfill
\begin{subfigure}[b]{0.48\textwidth}
  \centering
  \adjustbox{max width=\linewidth}{% 全局协调统一右绕 → 通行（子图 b）
\begin{tikzpicture}[scale=0.85, transform shape]
  % 道路
  \fill[bglight] (-0.5,-2.5) rectangle (8.5,2.5);
  \draw[gray, thick] (-0.5, 2.5) -- (8.5, 2.5);
  \draw[gray, thick] (-0.5,-2.5) -- (8.5,-2.5);
  \draw[dashed, gray!60] (-0.5, 0) -- (8.5, 0);
  \node[gray, font=\scriptsize, anchor=east] at (-0.6, 0) {$l{=}0$};

  % O_a（上方靠边）y=[0.5, 1.0]
  \fill[dangerred!70, draw=dangerred, thick] (3.2, 0.5) rectangle (4.8, 1.0);
  \node[font=\scriptsize, white] at (4.0, 0.75) {$O_a$};

  % O_b（下方靠边）y=[-1.0, -0.5]
  \fill[dangerred!70, draw=dangerred, thick] (3.2,-1.0) rectangle (4.8,-0.5);
  \node[font=\scriptsize, white] at (4.0,-0.75) {$O_b$};

  % 膨胀虚线 δ=0.4（统一右绕）
  \draw[dangerred!60, dashed, thick] (2.8, 0.1) rectangle (5.2, 1.4);
  \draw[dangerred!60, dashed, thick] (2.8,-1.4) rectangle (5.2,-0.1);

  % 决策标注（统一 R）
  \node[deepblue, font=\scriptsize\bfseries] at (5.5, 0.75) {R};
  \node[deepblue, font=\scriptsize\bfseries] at (5.5,-0.75) {R};

  % 统一右绕 → 上界压到 O_b 膨胀下缘 -1.4（offset -0.1）
  \draw[deeppink, very thick]
    (-0.5, 2.5) -- (2.6, 2.5) -- (2.6, -1.5) -- (5.4, -1.5) -- (5.4, 2.5) -- (8.5, 2.5);
  \node[deeppink, font=\scriptsize] at (7.5, 2.1) {$l^+(s)$};

  % 下界保持道路下界
  \draw[deepblue, very thick, dashed]
    (-0.5,-2.5) -- (8.5,-2.5);
  \node[deepblue, font=\scriptsize, left] at (-0.6,-2.5) {$l^-(s)$};

  % 通行带
  \fill[lightgreen, opacity=0.25] (-0.5,-2.5) rectangle (8.5,-1.5);

  % 通行带宽度
  \draw[lightgreen!60!black, thick, {Stealth}-{Stealth}] (0.5,-2.5) -- (0.5,-1.5);
  \node[lightgreen!60!black, font=\scriptsize, anchor=east] at (0.4,-2.0) {$W{=}1.0$m};

  % 自车轨迹（无 fill）
  \draw[deepblue, very thick, -{Stealth}]
    plot[smooth, tension=0.55] coordinates
      {(0.5, 0) (1.5,-0.5) (2.5,-1.8) (4.0,-2.0) (5.5,-1.8) (6.5,-0.5) (7.5, 0)};
  \node[deepblue, font=\scriptsize] at (4.0,-2.4) {规划轨迹};

  % PASS 判定
  \node[draw=lightgreen!80!black, fill=lightgreen!20, font=\scriptsize,
        rounded corners=3pt] at (4.0, 2.0)
    {统一右绕 $\Rightarrow$ \textbf{PASS}};
\end{tikzpicture}
}
  \caption{组级协同选向恢复通行带}
\end{subfigure}
\caption{本文处理的多障碍物空间选向冲突}
\caption*{逐个不可回退的选向会同时压缩路径上下界，而将组内方向联合成连续分割后可以直接保留有效间隙。}
\label{fig:q34_teaser}
\end{figure}

% ====== 2 相关工作与问题定义 ======
\section{相关工作与问题定义}
\label{sec:related}
\label{sec:problem}

\subsection{相关工作与研究缺口}

SLT 等时空表示上的联合搜索与优化可以直接协调几何和到达时序，但通常需要更换或扩展原解耦求解链\citep{mcnaughton2011lattice,yang2021_ia_planner,hu2023_hybrid_astar_slt,hu2025_dp_nmpc,qiao2025stjoint}，而迭代更新、候选配对和软代价方法虽然保留路径--速度分解，却需要多轮调用或额外候选生成\citep{liu2017speed,chen2019em,zhang2023_dp_rcg,han2026_contingency,pathspeedstructured2024}。混合整数规划能够统一离散避让方向与连续轨迹，但其计算和实现结构也不同于现有两阶段 QP\citep{kessler2023}。

Frenet 或时空安全走廊以显式可行域连接离散决策与连续优化\citep{liu2017sfc,yoon2024stcorridor,frenetcorridor2025}，并与关注动态预测和障碍物交互的局部规划研究形成互补\citep{schwarting2018planning,yangbo2023ped,jeong2021ml}。现有方法仍缺少一种适配 Apollo 等解耦规划器\citep{fan2018baidu,apollo2024}的统一接口，能够在不更换两阶段 QP 的条件下联合处理到达时刻、组级绕行及其速度时序约束。

% ====== 3 问题提出 ======
\subsection{解耦规划接口与失效机理}
\label{subsec:decouple_formalize}

路径--速度解耦先在纵向--横向（station--lateral, SL）空间规划横向路径 $l(s)$，再沿该路径在纵向--时间（station--time, ST）空间规划速度曲线 $s(t)$，因而将三维时空轨迹规划转化为两个顺序求解的二维子问题。通用求解主干在空间采样点 $s_j$ 上优化路径状态 $\mathbf x_p=(l_j,l'_j,l''_j)_{j=0}^{K}$，并在时间采样点 $t_r$ 上优化速度状态 $\mathbf x_v=(s_r,v_r,a_r)_{r=0}^{N}$，相应的分段加加速度 QP 可统一写为\citep{fan2018baidu,zhou2021pathqp}

\begin{align}
 \mathbf x_p^*={}&\arg\min_{\mathbf x_p}
 \frac{1}{2}\mathbf x_p^{\mathsf T}H_p\mathbf x_p+q_p^{\mathsf T}\mathbf x_p,
 &\mathrm{s.t.}\quad&A_p\mathbf x_p=b_p,\quad l_j^-\leq l_j\leq l_j^+, \label{eq:path_qp}\\
 \mathbf x_v^*={}&\arg\min_{\mathbf x_v}
 \frac{1}{2}\mathbf x_v^{\mathsf T}H_v\mathbf x_v+q_v^{\mathsf T}\mathbf x_v,
 &\mathrm{s.t.}\quad&A_v\mathbf x_v=b_v,\quad s_r^{lb}\leq s_r\leq s_r^{ub}, \notag\\
 &&&0\leq v_r\leq v_{max},\quad a_{min}\leq a_r\leq a_{max}. \label{eq:speed_qp}
\end{align}
其中 $H_p$、$H_v$ 编码平滑代价，$A_p$、$A_v$ 编码相邻采样点间的分段加加速度运动学等式，而路径边界 $\mathcal B_l=\{[l_j^-,l_j^+]\}$ 与 ST 位置边界 $\mathcal B_s=\{[s_r^{lb},s_r^{ub}]\}$ 则是两个凸 QP 接收环境约束的核心接口。MIKU 保留式\eqref{eq:path_qp}--\eqref{eq:speed_qp}的目标矩阵和运动学等式，并按下式重构两组边界。
\begin{equation}
 (\mathcal B_l,\mathcal H_s)=\Phi_{\mathrm{MIKU}}(\mathcal R,\mathcal X_{ego},\{O_i(t)\}),
 \qquad
 \mathcal H_t=\Phi_{\mathrm{time}}(l(s),\{O_i(t)\}),\quad
 \mathcal B_s=\Phi_{\mathrm{ST}}(l(s),\{O_i(t)\})\cap\Phi_{\mathcal H_t}(\mathcal H_t),
 \label{eq:miku_interface}
\end{equation}
其中 $\mathcal R$ 为道路与参考线信息，$\mathcal X_{ego}$ 为自车状态，$\mathcal H_s$、$\mathcal H_t$ 分别为空间与时间同伦，$\Phi_{\mathrm{ST}}$ 表示基础规划器的路径相关 ST 映射，$\Phi_{\mathcal H_t}$ 表示将时间同伦写入 ST 上下边界的跨阶段传递。式\eqref{eq:miku_interface}由此把环境预测、空间选向和速度时序汇入同一解耦求解链。

在 Frenet 坐标系下，第 $i$ 个障碍物的物理投影记为 $O_i=[s_i^{-},s_i^{+}]\times[l_i^{-},l_i^{+}]$，道路物理横向边界记为 $[l_{road}^{-},l_{road}^{+}]$。设自车宽度为 $W_{ego}$，路侧额外裕度为 $d_{road}$，则自车中心可行道路区间为
\begin{equation}
    [L^{-},L^{+}]=
    [l_{road}^{-}+W_{ego}/2+d_{road},\;l_{road}^{+}-W_{ego}/2-d_{road}].
    \label{eq:center_road}
\end{equation}
障碍物物理边界同样只外扩一次半车宽与障碍物额外裕度 $d_{buf,i}$，得到针对自车中心的禁行区间
\begin{equation}
    [u_i,v_i]=[l_i^{-}-W_{ego}/2-d_{buf,i},\;l_i^{+}+W_{ego}/2+d_{buf,i}].
    \label{eq:center_forbidden}
\end{equation}
因而 $L^{-}$、$L^{+}$、$u_i$、$v_i$ 均表示自车中心的可行或禁行边界，后续间隙宽度表示中心可放置区间的长度，其中车宽已包含在上述一次内缩与外扩之中。

为每个在截面 $s$ 处活跃的障碍物分配二值避让决策 $d_i\in\{L,R\}$，记两类集合为 $\mathcal L$ 和 $\mathcal R$，则中心通行带为
\begin{equation}
    l^{-}(\mathbf d)=\max\!\left(L^{-},\max_{i\in\mathcal L}v_i\right),\qquad
    l^{+}(\mathbf d)=\min\!\left(L^{+},\min_{i\in\mathcal R}u_i\right).
    \label{eq:center_band}
\end{equation}
其宽度 $W(\mathbf d)=l^{+}(\mathbf d)-l^{-}(\mathbf d)$，给定数值容差 $\varepsilon>0$ 后，$W(\mathbf d)\geq\varepsilon$ 即表示存在可用的自车中心区间，因为车宽已经通过式\eqref{eq:center_road}和式\eqref{eq:center_forbidden}纳入边界。

由于真实障碍物随时间运动，本文用自车到达纵向位置 $s$ 的预估时刻 $\tau(s)$ 查询 $O_i(\tau(s))$，并按式\eqref{eq:center_forbidden}同步更新禁行边界。对包含 $k$ 个障碍物的纵向连通分量 $C$，记障碍物 $i$ 的代表位置为 $\bar s_i$，并在 $\bar t_i=\tau(\bar s_i)$ 处查询 $[u_i(\bar t_i),v_i(\bar t_i)]$，则组级选向问题为
\begin{equation}
    \mathbf{d}_{C}^{*} = \arg\max_{\mathbf{d}_{C}\in\{L,R\}^{k}}
    \left[\min\!\left(L^{+},\min_{i\in\mathcal R_C}u_i(\bar t_i)\right)
    -\max\!\left(L^{-},\max_{i\in\mathcal L_C}v_i(\bar t_i)\right)\right],
    \label{eq:prob_joint}
\end{equation}
其中 $\mathbf d_C=(d_1,\ldots,d_k)$ 为组内绕行方向，$\mathcal L_C$、$\mathcal R_C$ 为对应分组。逐项枚举需要检查 $2^k$ 种组合，而第\ref{sec:method_band}节将证明式\eqref{eq:prob_joint}存在连续分割最优解，并以 $O(k\log k)$ 的排序与扫描求得最大中心间隙。

\label{subsec:union_projection}

路径子问题忽略时间时，通常以预测时域 $[0,T_{\text{pred}}]$ 上的全时刻并集近似动态占据
\begin{equation}
    \bar{O}_i = \bigcup_{t \in [0,\,T_{\text{pred}}]} O_i(t),
    \label{eq:prob_union}
\end{equation}
式~\eqref{eq:prob_union} 将占据的时间条件压缩为曾经占据的区域，而只取 $t{=}0$ 快照或暂不让动态障碍物参与路径边界，也同样无法沿自车到达时间查询预测轨迹。图~\ref{fig:dynamic_exclusion} 中的行人在自车到达相同纵向位置前已让出通道，但全时刻并集仍将其扫掠区域视为同时占据，因而把可行的时序窗口误判为阻塞。

\begin{figure}[htbp]
\centering
\begin{subfigure}[b]{0.30\textwidth}
  \centering
  \adjustbox{max width=\linewidth}{% 动态障碍物排除 t=0 时刻（子图 a）
\begin{tikzpicture}[scale=0.95, transform shape]
  \fill[bglight] (0,-1.3) rectangle (6.8,1.8);
  \draw[deepblue, thick] (0,-1.3) -- (6.8,-1.3);
  \draw[deepblue, thick] (0, 1.8) -- (6.8, 1.8);
  \draw[deepblue, dashed, thin] (0,0.3) -- (6.8,0.3);

  \fill[egopt] (0.5,0) rectangle (1.35,0.6);
  \node[white, font=\scriptsize\bfseries] at (0.925,0.3) {ADC};
  \draw[-Stealth, thick, deepblue] (1.4,0.3) -- (2.5,0.3);
  \node[font=\scriptsize, deepblue] at (1.95,0.55) {8 m/s};

  \fill[obsdyn] (5.7,-0.2) rectangle (6.1,0.3);
  \draw[-Stealth, thick, dangerred] (5.9,0.35) -- (5.9,1.4);
  \node[font=\scriptsize, dangerred!70!black] at (6.3,1.2) {1.2 m/s};

  \draw[<->, thin, gray] (1.4,-0.75) -- (5.7,-0.75);
  \node[font=\scriptsize, gray] at (3.55,-0.65) {$s_{\text{rel}}=10$ m};
\end{tikzpicture}
}
  \caption{$t{=}0$ 行人未进入路径}
\end{subfigure}\hfill
\begin{subfigure}[b]{0.30\textwidth}
  \centering
  \adjustbox{max width=\linewidth}{% 动态障碍物排除 t=0.625s（子图 b）
\begin{tikzpicture}[scale=0.95, transform shape]
  \fill[bglight] (0,-1.3) rectangle (6.8,1.8);
  \draw[deepblue, thick] (0,-1.3) -- (6.8,-1.3);
  \draw[deepblue, thick] (0, 1.8) -- (6.8, 1.8);
  \draw[deepblue, dashed, thin] (0,0.3) -- (6.8,0.3);

  \fill[egopt] (3.4,0) rectangle (4.25,0.6);
  \node[white, font=\scriptsize\bfseries] at (3.825,0.3) {ADC};
  \draw[-Stealth, thick, deepblue] (4.3,0.3) -- (5.4,0.3);

  \fill[obsdyn] (5.7,0.4) rectangle (6.1,0.85);
  \draw[-Stealth, thick, dangerred] (5.9,0.9) -- (5.9,1.5);

  \draw[dangerred, very thick, dashed] (5.55,-0.2) rectangle (6.25,1.0);
\end{tikzpicture}
}
  \caption{中间时刻行人半入路径}
\end{subfigure}\hfill
\begin{subfigure}[b]{0.30\textwidth}
  \centering
  \adjustbox{max width=\linewidth}{% 动态障碍物排除 t=1.25s（子图 c）
\begin{tikzpicture}[scale=0.95, transform shape]
  \fill[bglight] (0,-1.3) rectangle (6.8,1.8);
  \draw[deepblue, thick] (0,-1.3) -- (6.8,-1.3);
  \draw[deepblue, thick] (0, 1.8) -- (6.8, 1.8);
  \draw[deepblue, dashed, thin] (0,0.3) -- (6.8,0.3);

  \fill[egopt, opacity=0.3] (0.5,0) rectangle (1.35,0.6);
  \node[gray, font=\scriptsize] at (0.925,0.3) {$t=0$};

  \fill[egopt] (5.7,0) rectangle (6.55,0.6);
  \node[white, font=\scriptsize\bfseries] at (6.125,0.3) {ADC};

  \fill[obsdyn] (5.7,1.3) rectangle (6.1,1.75);
  \node[font=\scriptsize, dangerred!70!black, above] at (5.9,1.8) {已穿过};

  \fill[bandok] (1.5,0.05) rectangle (5.6,0.55);
  \node[font=\scriptsize, lightgreen!50!black] at (3.5,0.3) {路径已打开};
\end{tikzpicture}
}
  \caption{自车到达时通行带打开}
\end{subfigure}
\caption{横穿行人场景下的时序约束缺失}
\caption*{全时刻并集投影把行人扫过的车道并入固定足迹，而沿到达时间 $\tau(s)$ 的时变投影只查询自车实际到达时刻的占用，此时通行带已经打开。}
\label{fig:dynamic_exclusion}
\end{figure}

时间无关投影使路径阶段无法区分当前占据与到达时占据，而逐障碍物贪心则只检查 $2^k$ 种方向组合中的一种，当不可回退的局部选择同时压缩上下边界时，路径子问题可能因 $l^{+}(s)<l^{-}(s)$ 而进入停车降级。本文由此求解两个相互关联的问题，即沿 $\tau(s)$ 恢复动态占据的时间条件，并在同一纵向连通分量内联合选择绕行方向，从而把真实阻塞与时间压缩或局部选向造成的虚假不可行区分开来。

% ====== 3 MIKU 算法（方法 + 理论） ======
\section{MIKU 多障碍物时空同伦约束统一方法}
\label{sec:method}

在第\ref{sec:problem}节的中心边界语义下，MIKU 先估计 $\tau(s)$ 并查询时变 SL 投影，再通过 $O(k\log k)$ 最大间隙扫描和 Top-$K$ 空间同伦生成连续路径边界，路径求解后沿 $l(s)$ 构造预测占用管与安全时间窗，并以分层时间同伦图选择全局因果序列写入 ST 上下边界。图~\ref{fig:miku_method_flow}概括从障碍物预测到路径边界和 ST 位置边界的约束构造与传递关系，算法~\ref{alg:miku}给出相应步骤，其中离散同伦被转换为两组边界，连续凸 QP 则据此生成平滑轨迹。

\begin{figure}[htbp]
\centering
\adjustbox{max width=\textwidth}{% 三区/四区稿专用：平台无关的 MIKU 算法主链
\begin{tikzpicture}[
  node distance=8mm and 7mm,
  input/.style={rectangle,rounded corners=2pt,draw=deepblue,fill=inputyellow,
    minimum width=2.6cm,minimum height=.9cm,align=center,font=\footnotesize},
  miku/.style={rectangle,rounded corners=3pt,draw=deeppink,very thick,
    fill=improvedpink!55,minimum width=3.05cm,minimum height=1.05cm,
    align=center,font=\footnotesize},
  solver/.style={rectangle,rounded corners=3pt,draw=deepblue,
    fill=lightblue!50,minimum width=2.65cm,minimum height=1.0cm,
    align=center,font=\footnotesize},
  output/.style={rectangle,rounded corners=2pt,draw=deepblue,fill=outputpink,
    minimum width=2.7cm,minimum height=.9cm,align=center,font=\footnotesize},
  algflow/.style={-Stealth,very thick,deeppink},
  solveflow/.style={-Stealth,thick,deepblue},
  note/.style={font=\scriptsize,align=center,fill=white,inner sep=1pt}
]

\node[input] (input) {道路几何、自车状态\\多障碍物预测};
\node[miku, right=of input] (occupancy) {到达时间查询\\威胁裕度与鲁棒占据};
\node[miku, right=of occupancy] (spatial) {连续分割扫描\\Top-K 空间同伦};
\node[solver, right=of spatial] (pathqp) {路径凸 QP\\求解 $l(s)$};

\node[miku, below=of spatial] (temporal) {安全时间窗\\多冲突时间同伦};
\node[solver, right=of temporal] (speedqp) {双向 ST 边界\\速度凸 QP};
\node[output, right=of speedqp] (trajectory) {时空轨迹\\同伦语义标签};

\draw[algflow] (input) -- (occupancy);
\draw[algflow] (occupancy) -- (spatial);
\draw[solveflow] (spatial) -- node[note,above]{中心路径边界} (pathqp);
\draw[solveflow] (pathqp) |- node[note,pos=.56,right, xshift=2pt]{$l(s)$ 与冲突区} (temporal);
\draw[algflow] (temporal) -- (speedqp);
\draw[solveflow] (speedqp) -- (trajectory);
% 两条反馈边放在图外侧，并将说明文字置于空白区域，避免穿过节点。
\draw[algflow,rounded corners]
  (speedqp.south) -- ++(0,-.55) -- ++(-7.0,0) -- ++(0,1.05)
  node[pos=.18,below, note] {速度解更新 $\tau(s)$} -| (occupancy.south);
\draw[solveflow,rounded corners]
  (trajectory.north) -- ++(0,.45) -- ++(-1.0,0)
  node[pos=.45,above, note] {滚动同伦保持} -| (temporal.east);
\end{tikzpicture}
}
\caption{MIKU 的约束构造与传递流程}
\caption*{红色模块表示 MIKU 的候选与约束构造，蓝色模块表示路径和速度连续求解器。算法先以空间同伦生成路径边界，再以时间同伦将通行时序写入 ST，速度解可在需要时更新下一轮到达时间估计。}
\label{fig:miku_method_flow}
\end{figure}

\subsection{交互感知威胁度与差异化安全裕度}
\label{sec:method_threat}

统一缓冲在窄路两侧同时受限时容易过度挤压净间距，因此 MIKU 以碰撞紧迫性、横向重叠、相对运动、目标类型和局部交互密度构造归一化威胁度 $\Theta_i$，并据此为不同障碍物生成差异化安全裕度 $d_{buf,i}$。

\begin{equation}
    \Theta_i = w_1 f_{TTC}(i) + w_2 f_{overlap}(i) + w_3 f_{vel}(i) + w_4 f_{type}(i) + w_5 f_{inter}(i),
    \label{eq:threat_score}
\end{equation}

其中各因子归一化至 $[0,1]$，权重之和为一。本文取 $(w_1,w_2,w_3,w_4,w_5)=(0.30,0.20,0.15,0.10,0.25)$，依次体现碰撞紧迫性、横向重叠、相对速度、目标类型和局部密度在边界扩张中的优先级，其权重敏感性在第\ref{subsec:exp_sensitivity}节分析。

碰撞时间（Time To Collision, TTC）因子衡量纵向冲突的紧迫性，当 $s_i>s_{ego}$ 且 $\dot{s}_{ego}>\dot{s}_i$ 时取 $\mathrm{TTC}_i=(s_i-s_{ego})/(\dot{s}_{ego}-\dot{s}_i)$，其余不构成追赶冲突的情形令 $\mathrm{TTC}_i=+\infty$。分段线性映射在 $\mathrm{TTC}_i\leq T_{crit}$ 时令 $f_{TTC}=1$，在 $\mathrm{TTC}_i\geq T_{max}$ 时令 $f_{TTC}=0$，中间区间线性递减，其中依据碰撞时间分布分析\citep{minderhoud2001extended}取 $T_{crit}=2.0$\,s，并以覆盖速度规划时间域的 $T_{max}=7.0$\,s 为上限。

其余四个因子从横向重叠、相对速度、目标类型与局部密度补充威胁度。横向重叠率 $f_{overlap}(i)=\max\!\big(0,\min(l_i^{+},l_{ego}^{+})-\max(l_i^{-},l_{ego}^{-})\big)/W_{ego}$ 度量障碍物与自车走廊的重叠比例，而相对速度因子 $f_{vel}(i)=\sigma(v_{rel,i}/v_{max})$ 采用 $\sigma(x)=1/(1+e^{-5x})$ 的 Sigmoid 映射，其中 $v_{rel,i}=\dot{s}_{ego}-\dot{s}_i$，$v_{max}$ 为归一化尺度。类型因子按保护优先级为弱势道路使用者、机动车、未知可动目标、静态障碍物和小尺寸标识物分别取 $1.0$、$0.7$、$0.5$、$0.3$ 和 $0.15$，交互密度因子 $f_{inter}(i)=\frac{1}{N-1}\sum_{j\neq i}\mathbf{1}[d_{ij}<d_{cluster}](d_{cluster}-d_{ij})/d_{cluster}$ 则以 $d_{cluster}=10$\,m 加权统计邻近障碍物，且在 $N=1$ 时取零。

威胁度确定后，额外缓冲在区间 $[d_{buf,min},d_{buf,max}]$ 上线性插值。

\begin{equation}
    d_{buf,i} = d_{buf,min} + (d_{buf,max} - d_{buf,min})\,\Theta_i,
    \label{eq:dynamic_delta}
\end{equation}

额外缓冲区间取 $[0.10,0.40]$\,m。半车宽 $W_{ego}/2$ 在式\eqref{eq:center_forbidden}中始终保留，式\eqref{eq:dynamic_delta}仅改变障碍物外侧的附加空间，代入式\eqref{eq:center_forbidden}得

\begin{equation}
    u_i=l_i^{-}-W_{ego}/2-d_{buf,i},\qquad
    v_i=l_i^{+}+W_{ego}/2+d_{buf,i},
    \label{eq:dynamic_uv}
\end{equation}

其中 $u_i$、$v_i$ 已是自车中心约束。$d_{buf,i}$ 只改变边界取值，不改变第\ref{sec:method_band}节的排序与扫描结构。图~\ref{fig:threat_radar}展示该参数化效果。

\begin{figure}[htbp]
\centering
\begin{subfigure}[b]{0.42\textwidth}
  \centering
  \adjustbox{max width=\linewidth}{% 高威胁障碍物雷达图（子图 a）
\begin{tikzpicture}[scale=0.85, transform shape]
  \foreach \r in {1, 2, 3} {
    \draw[gray!40, thin] (90:\r) \foreach \a in {162, 234, 306, 378} { -- (\a:\r) } -- cycle;
  }
  \foreach \a/\lbl in {90/$f_{TTC}$, 162/$f_{overlap}$, 234/$f_{vel}$, 306/$f_{type}$, 378/$f_{inter}$} {
    \draw[gray!60, thin] (0,0) -- (\a:3.3);
    \node[font=\footnotesize] at (\a:3.7) {\lbl};
  }

  \fill[dangerred!30, draw=dangerred, thick]
    (90:2.7) -- (162:2.4) -- (234:2.55) -- (306:2.1) -- (378:1.8) -- cycle;
  \foreach \a/\r in {90/2.7, 162/2.4, 234/2.55, 306/2.1, 378/1.8} {
    \fill[dangerred] (\a:\r) circle (2.5pt);
  }

  \node[draw=dangerred, fill=dangerred!15, rounded corners, font=\small\bfseries] at (0, -4) {$\Theta = 0.82$ (CRITICAL)};

  \node[gray, font=\tiny] at (-3.8, 0) {0};
  \node[gray, font=\tiny] at (-3.8, 1.0) {0.33};
  \node[gray, font=\tiny] at (-3.8, 2.0) {0.67};
  \node[gray, font=\tiny] at (-3.8, 3.0) {1.0};
\end{tikzpicture}
}
  \caption{高威胁横穿行人}
\end{subfigure}\hfill
\begin{subfigure}[b]{0.42\textwidth}
  \centering
  \adjustbox{max width=\linewidth}{% 低威胁障碍物雷达图（子图 b）
\begin{tikzpicture}[scale=0.85, transform shape]
  \foreach \r in {1, 2, 3} {
    \draw[gray!40, thin] (90:\r) \foreach \a in {162, 234, 306, 378} { -- (\a:\r) } -- cycle;
  }
  \foreach \a/\lbl in {90/$f_{TTC}$, 162/$f_{overlap}$, 234/$f_{vel}$, 306/$f_{type}$, 378/$f_{inter}$} {
    \draw[gray!60, thin] (0,0) -- (\a:3.3);
    \node[font=\footnotesize] at (\a:3.7) {\lbl};
  }

  \fill[lightgreen!40, draw=lightgreen!70!black, thick]
    (90:0.3) -- (162:0.6) -- (234:0.3) -- (306:0.9) -- (378:0.45) -- cycle;
  \foreach \a/\r in {90/0.3, 162/0.6, 234/0.3, 306/0.9, 378/0.45} {
    \fill[lightgreen!70!black] (\a:\r) circle (2.5pt);
  }

  \node[draw=lightgreen!70!black, fill=lightgreen!15, rounded corners, font=\small\bfseries] at (0, -4) {$\Theta = 0.18$ (MINOR)};
\end{tikzpicture}
}
  \caption{低威胁远离静态目标}
\end{subfigure}
\caption{多因子威胁度画像与差异化安全裕度}
\caption*{五个归一化因子合成威胁度 $\Theta$，使高威胁横穿行人获得接近上限的安全裕度，而远离的静态目标采用接近下限的额外缓冲。}
\label{fig:threat_radar}
\end{figure}

\subsection{扫描线分组与组级最大间隙通行带}
\label{sec:method_band}

确定安全边界后，MIKU 将单截面决策改写为区间排序后的间隙选择，并按纵向重叠关系分组，使彼此无关的障碍物独立求解，同组障碍物则联合决策，从而避免局部选向分别压低上界和抬高下界所造成的虚假不可行。

当 $[s_i^{-},s_i^{+}]\cap[s_j^{-},s_j^{+}]\neq\varnothing$ 时，障碍物 $O_i$ 与 $O_j$ 纵向重叠，对该关系取传递闭包后可将 $n$ 个障碍物划为 $m$ 个连通分量 $C_1,\ldots,C_m$。分组按 $s_i^{-}$ 升序扫描，并维护当前分量的最大 $s^{+}$，若下一障碍物的 $s^{-}$ 不超过该值则并入当前分量，否则另起分量，因此排序与扫描的总复杂度为 $O(n\log n)$。

连通分量中的成员可以在不同截面依次活跃。MIKU 以组级决策统一障碍物的绕行侧，写回路径边界时则在每个 $s$ 处仅由当前活跃障碍物更新上下界。该方向一致性约束避免链式分量中的通行带频繁换侧，使离散边界保持连续、稳定。记分量 $C_j$ 在各活跃截面上生成的最小中心带宽度为 $\hat W(C_j)$，则整条离散路径的最小宽度为

\begin{equation}
    \hat{W} = \min_{j=1,\ldots,m} \hat{W}(C_j).
    \label{eq:global_width}
\end{equation}

在单个截面内，最优指在 $2^k$ 个方向组合中最大化式\eqref{eq:center_band}的中心带宽度。下述连续分割引理与最大间隙定理构成组级决策的理论基础。

\begin{lemma}[最优分割的连续性]
\label{lem:continuous_partition}
将障碍物按 $u_i$ 值升序排列后，存在最优解 $\mathbf{d}^{*}$，使得分割在排序中是连续的，即存在 $p \in \{0, 1, \ldots, k\}$ 满足
\begin{equation}
    d_{(i)}^{*} = L,\ \forall\, i \leq p, \qquad d_{(i)}^{*} = R,\ \forall\, i > p.
    \label{eq:continuous_partition}
\end{equation}
\end{lemma}

证明。两侧有一为空时取 $p=k$ 或 $p=0$。两侧均非空时令 $\beta=\min_{i\in\mathcal R}u_i$，将所有 $u_i\geq\beta$ 且原属 $\mathcal L$ 的障碍物改入 $\mathcal R$。由于新成员的 $u_i$ 不低于 $\beta$，式\eqref{eq:center_band}的上界不变，而 $\mathcal L$ 减少元素后下界不增，因而宽度不减。对并列 $u_i$ 用 $v_i$ 降序作二级排序，即可得到式\eqref{eq:continuous_partition}的连续分割最优解。$\hfill\Box$

引理 \ref{lem:continuous_partition} 将搜索空间由 $2^k$ 个任意分配降为 $k+1$ 个候选分割点 $p=0,1,\ldots,k$，其中两个端点表示全部障碍物位于通行带同侧，中间点表示从相邻障碍物之间通过。当 $u_i$ 相等时，以 $v_i$ 降序作为二级排序键即可保证排序唯一，而差异化的 $d_{buf,i}$ 不影响结论。

与中心道路区间 $[L^-,L^+]$ 完全不相交的禁行区间不会收紧当前截面，先从活跃列表中删除。其余区间允许越过某一侧道路边界，并由下式中的 $\min$ 和 $\max$ 自动完成道路裁剪。对分割点 $p$，由于按 $u_i$ 排序不能保证 $v_i$ 单调，定义包含道路下界的前缀最大边界
\begin{equation}
    V_0=L^{-},\qquad V_p=\max\!\left(L^{-},\max_{1\leq q\leq p}v_{(q)}\right),
    \label{eq:prefix_v}
\end{equation}
并定义候选上界 $U_p$ 与 $k+1$ 个中心间隙
\begin{equation}
    U_p=\begin{cases}
        \min(L^{+},u_{(p+1)}), & p=0,\ldots,k-1,\\
        L^{+}, & p=k,
    \end{cases}
    \qquad g_p=U_p-V_p,
    \label{eq:gap_def}
\end{equation}
其中 $V_p$ 必须使用前 $p$ 个 $v$ 的前缀最大值，因为相邻的 $v_{(p)}$ 会在禁行区间交错或嵌套时高估中间间隙，而已包含半车宽和道路裕度的 $L^-$ 与 $L^+$ 使 $g_p$ 可以直接表示自车中心可行宽度。

\begin{theorem}[最大间隙最优性]
\label{thm:max_gap}
对任一分割点 $p$，式\eqref{eq:center_band}的中心通行带下界为 $V_p$、上界为 $U_p$，故宽度 $W(p)=g_p$，且
\begin{equation}
    W^{*} = \max_{p = 0, 1, \ldots, k} g_p.
    \label{eq:max_gap_opt}
\end{equation}
\end{theorem}

证明。对分割 $p$，前 $p$ 个障碍物施加下界，式\eqref{eq:center_band}给出 $l^-(p)=V_p$，而升序排列的后 $k-p$ 个障碍物以 $u_{(p+1)}$ 形成最紧上界，故 $l^+(p)=U_p$ 且 $W(p)=g_p$。由引理\ref{lem:continuous_partition}可知至少存在一个连续分割最优解，因此遍历 $p=0,\ldots,k$ 后取最大 $g_p$ 即与 $2^k$ 组合枚举的最优宽度相同。$\hfill\Box$

定理\ref{thm:max_gap}中 $g_p<\varepsilon$ 表示该中心带不可用，仅当 $\max_p g_p<\varepsilon$ 时才判定截面阻塞。由于半车宽已进入式\eqref{eq:center_road}和式\eqref{eq:center_forbidden}，$g_p$ 直接用于中心区间的可用性判定，而排序、筛选与扫描使单截面组内求解的复杂度为 $O(k\log k)$。

单组最宽间隙并不必然组成纵向上最连续的整体路径。因此，MIKU 为每个连通分量 $C_j$ 保留宽度排名前 $Q$ 的正宽候选带 $B_{jq}=[a_{jq},b_{jq}]$，其中心为 $c_{jq}=(a_{jq}+b_{jq})/2$、宽度为 $g_{jq}=b_{jq}-a_{jq}$，实验取 $Q=3$。将纵向相邻分量视为分层图，通过动态规划选择空间同伦序列

\begin{equation}
 (q_1^*,\ldots,q_m^*)=arg\min_{q_1,\ldots,q_m}
 \left[-\sum_{j=1}^{m}g_{jq_j}
 +\lambda_s\left(|c_{1q_1}-l_0|+
 \sum_{j=2}^{m}|c_{jq_j}-c_{j-1,q_{j-1}}|\right)\right].
 \label{eq:spatial_homotopy}
\end{equation}
式\eqref{eq:spatial_homotopy}以节点代价奖励宽通道，并以边代价抑制相邻分量的带中心跳变。定义 $D_{jq}$ 为到达节点 $(j,q)$ 的最小累积代价，则递推关系为
\begin{equation}
 D_{jq}=-g_{jq}+\min_{r=1,\ldots,Q}
 \left[D_{j-1,r}+\lambda_s|c_{jq}-c_{j-1,r}|\right],
 \label{eq:spatial_dp}
\end{equation}
初层以 $D_{1q}=-g_{1q}+\lambda_s|c_{1q}-l_0|$ 初始化，末层取最小值并回溯即得式\eqref{eq:spatial_homotopy}在前 $K$ 个候选带上的全局最优序列。该层级保留定理\ref{thm:max_gap}给出的精确候选结构，同时避免把各组局部最宽解机械串联成频繁换侧的路径。

\begin{figure}[htbp]
\centering
\adjustbox{max width=0.8\textwidth}{% 引理1+定理1 几何图证（1D 横向数轴）
\begin{tikzpicture}[scale=1.0, transform shape]

% 障碍物位置设计: g₂ 为最大间隙
% O₁@1.5, O₂@3.0, O₃@6.5, O₄@8.0 (中心)
% 半宽 0.35, δ 扩展 0.55
% l_road⁻=0, l_road⁺=9.5

% === 上图: 排序后的障碍物横向投影 ===

% 道路区域淡色背景（先画，避免盖住坐标轴）
\fill[bglight] (0,-0.7) rectangle (9.5,0.7);

% l 轴
\draw[->, thick] (-0.5,0) -- (10.2,0) node[right,font=\footnotesize] {$l$};

% 道路边界（垂直于 l 轴的竖线）
\draw[deepblue, very thick] (0,-0.9) -- (0,0.9);
\node[deepblue, font=\scriptsize, below] at (0,-0.95) {$l_{road}^-$};
\draw[deepblue, very thick] (9.5,-0.9) -- (9.5,0.9);
\node[deepblue, font=\scriptsize, below] at (9.5,-0.95) {$l_{road}^+$};

% 4 个障碍物（按 u_i 排序）
\fill[obsstatic] (1.15,-0.35) rectangle (1.85,0.35);
\node[white, font=\scriptsize\bfseries] at (1.5,0) {$O_1$};

\fill[obsstatic] (2.65,-0.35) rectangle (3.35,0.35);
\node[white, font=\scriptsize\bfseries] at (3.0,0) {$O_2$};

\fill[obsstatic] (6.15,-0.35) rectangle (6.85,0.35);
\node[white, font=\scriptsize\bfseries] at (6.5,0) {$O_3$};

\fill[obsstatic] (7.65,-0.35) rectangle (8.35,0.35);
\node[white, font=\scriptsize\bfseries] at (8.0,0) {$O_4$};

% u_i / v_i 膨胀区
\draw[obsbuf, dashed] (0.95,-0.55) rectangle (2.05,0.55);
\draw[obsbuf, dashed] (2.45,-0.55) rectangle (3.55,0.55);
\draw[obsbuf, dashed] (5.95,-0.55) rectangle (7.05,0.55);
\draw[obsbuf, dashed] (7.45,-0.55) rectangle (8.55,0.55);

% u_i / v_i 标注（选择性标注避免拥挤）
\node[deepblue, font=\tiny] at (0.95,0.75) {$u_1$};
\node[deeppink, font=\tiny] at (2.05,0.75) {$v_1$};
\node[deepblue, font=\tiny] at (2.45,0.75) {$u_2$};
\node[deeppink, font=\tiny] at (3.55,0.75) {$v_2$};
\node[deepblue, font=\tiny] at (5.95,0.75) {$u_3$};
\node[deeppink, font=\tiny] at (7.05,0.75) {$v_3$};
\node[deepblue, font=\tiny] at (7.45,0.75) {$u_4$};
\node[deeppink, font=\tiny] at (8.55,0.75) {$v_4$};

% 5 个间隙标注（在下方）
% g₀ = u₁ - l_road⁻ = 0.95
\draw[<->, thick, lightgreen!60!black] (0,-1.2) -- (0.95,-1.2);
\node[font=\scriptsize, lightgreen!60!black] at (0.475,-1.5) {$g_0$};

% g₁ = u₂ - v₁ = 0.40
\draw[<->, thick, lightgreen!60!black] (2.05,-1.2) -- (2.45,-1.2);
\node[font=\scriptsize, lightgreen!60!black] at (2.25,-1.5) {$g_1$};

% g₂ = u₃ - v₂ = 2.40 ← 最大间隙
\draw[<->, very thick, deeppink] (3.55,-1.2) -- (5.95,-1.2);
\node[font=\small\bfseries, deeppink] at (4.75,-1.5) {$g_2^\star$};

% g₃ = u₄ - v₃ = 0.40
\draw[<->, thick, lightgreen!60!black] (7.05,-1.2) -- (7.45,-1.2);
\node[font=\scriptsize, lightgreen!60!black] at (7.25,-1.5) {$g_3$};

% g₄ = l_road⁺ - v₄ = 0.95
\draw[<->, thick, lightgreen!60!black] (8.55,-1.2) -- (9.5,-1.2);
\node[font=\scriptsize, lightgreen!60!black] at (9.025,-1.5) {$g_4$};

% 最优分割线
\draw[deeppink, thick, dashed] (4.75,-0.9) -- (4.75,0.9);

\node[font=\scriptsize] at (4.75,-1.9) {$k{=}4$ 个障碍物，$k{+}1{=}5$ 个候选间隙，$g_2^\star$ 最大};

% === 下图: 决策向量 ===
\begin{scope}[shift={(0,-3.8)}]

  % R 群背景
  \fill[lightgreen!20] (0,-0.6) rectangle (4.75,0.6);
  % L 群背景
  \fill[improvedpink!25] (4.75,-0.6) rectangle (9.5,0.6);

  % 4 个障碍物
  \fill[obsstatic] (1.15,-0.35) rectangle (1.85,0.35);
  \node[white, font=\scriptsize\bfseries] at (1.5,0) {$O_1$};
  \node[deepblue, font=\scriptsize\bfseries] at (1.5,-0.65) {$d_1{=}R$};

  \fill[obsstatic] (2.65,-0.35) rectangle (3.35,0.35);
  \node[white, font=\scriptsize\bfseries] at (3.0,0) {$O_2$};
  \node[deepblue, font=\scriptsize\bfseries] at (3.0,-0.65) {$d_2{=}R$};

  \fill[obsstatic] (6.15,-0.35) rectangle (6.85,0.35);
  \node[white, font=\scriptsize\bfseries] at (6.5,0) {$O_3$};
  \node[deeppink, font=\scriptsize\bfseries] at (6.5,-0.65) {$d_3{=}L$};

  \fill[obsstatic] (7.65,-0.35) rectangle (8.35,0.35);
  \node[white, font=\scriptsize\bfseries] at (8.0,0) {$O_4$};
  \node[deeppink, font=\scriptsize\bfseries] at (8.0,-0.65) {$d_4{=}L$};

  % 分割线
  \draw[deeppink, thick, dashed] (4.75,-0.8) -- (4.75,0.8);
  \node[deeppink, font=\scriptsize] at (4.75,1.0) {$p^\star{=}2$};

  % 群标注
  \node[deepblue, font=\scriptsize] at (2.25,-1.1) {$\mathcal{R}{=}\{O_1,O_2\}$};
  \node[deeppink, font=\scriptsize] at (7.25,-1.1) {$\mathcal{L}{=}\{O_3,O_4\}$};

  % 连续性说明
  \node[font=\scriptsize] at (4.75,-1.5) {$\mathbf{d}^\star{=}(R,R,L,L)$，排序中连续，无交错};
\end{scope}

\end{tikzpicture}
}
\caption{按 $u_i$ 排序后的候选间隙与最优连续分割}
\caption*{$k{=}4$ 个障碍物在排序后产生 $k{+}1{=}5$ 个候选间隙，其中最大间隙 $g_2^{\star}$ 对应连续无交错的最优分割点 $p^{\star}{=}2$。}
\label{fig:max_gap}
\end{figure}

\subsection{到达时间驱动的时变投影}
\label{sec:method_corridor}

时变投影以到达时间函数 $\tau(s)$ 查询自车抵达纵向位置 $s$ 时的障碍物占据，本文机制场景由当前自车状态的二阶运动学模型计算该函数，滚动规划时也可采用上周期速度曲线。以匀加速运动学估计时

\begin{equation}
    \tau(s) = \frac{-v_{ego} + \sqrt{v_{ego}^{2} + 2\,a_{ego}\,(s - s_{ego})}}{a_{ego}},
    \label{eq:arrival_time}
\end{equation}

其中 $s_{ego}$、$v_{ego}$、$a_{ego}$ 来自当前状态估计，当 $|a_{ego}|<10^{-3}$ 时采用 $\tau(s)=(s-s_{ego})/\max(v_{ego},10^{-3})$，当判别式小于零时返回超出规划时域的大时刻标记。若使用上周期单调速度曲线，则对首次满足 $s_{prev}(t)\geq s$ 的时刻作一维插值，并在每个规划周期更新。

本文场景以常速度模型 $s_i(t)=s_{i,0}+v_{s,i}t$ 和 $l_i(t)=l_{i,0}+v_{l,i}t$ 给出障碍物预测，并结合几何尺寸得到时变 SL 边界 $[l_i^{-}(\tau(s)),l_i^{+}(\tau(s))]$，若该时刻的纵向占据不覆盖截面 $s$，则将障碍物标记为不活跃。该查询形式也可直接接收上游预测器输出的离散轨迹及其插值结果，静态障碍物则保留当前观测位置。

为使时变通行带在有界预测误差下仍保持几何安全性，MIKU 将障碍物中心预测扩张为不确定性占用管。设初始位置误差界为 $(\epsilon_{s0,i},\epsilon_{l0,i})$，速度误差界为 $(\epsilon_{vs,i},\epsilon_{vl,i})$，则时刻 $t$ 的纵横向扩张量为
\begin{equation}
 \epsilon_{s,i}(t)=\epsilon_{s0,i}+t\epsilon_{vs,i},\qquad
 \epsilon_{l,i}(t)=\epsilon_{l0,i}+t\epsilon_{vl,i}.
 \label{eq:prediction_tube}
\end{equation}
障碍物物理外形与式\eqref{eq:prediction_tube}、自车外形依次作 Minkowski 扩张，得到进入 ST 映射的预测占用管。因此，MIKU 使用的时间窗边界来自显式误差界下的障碍物占用，而不是由名义到达时间代替占用预测。

在每个路径离散位置 $s$，仅当障碍物在 $\tau(s)$ 时刻的纵向占据覆盖该截面时，才用其横向边界构造式\eqref{eq:center_forbidden}。组级分割在代表纵向位置处联合选择，并在分量内保持各障碍物的绕行侧不变，而写回路径边界时的 $u_i(\tau(s))$、$v_i(\tau(s))$ 与活跃集合仍逐截面重算，从而将代表区间上的精确最大间隙稳定地延伸到各活跃截面。

路径边界因此成为随 $\tau(s)$ 更新的动态通行带，已经让出的横向空间可以在对应截面释放，而通行带成立所需的到达条件则由速度规划处理。

当采用交替到达时间更新时，首次路径和速度求解得到的 $s(t)$ 用于重新构造到达时间估计 $\tau(s)$，再以阻尼系数 $\gamma$ 将该估计与上一轮的 $\tau(s)$ 融合并重新求解约束问题。完整 MIKU 最多执行两轮，若更新后的到达时间变化不超过收敛阈值或当前候选已经满足规划判据，则保留上一轮与当前轮中规划评分较高的结果，关闭该更新即得到消融变体 A7。

\subsection{多障碍物时间同伦与 ST 上下边界传递}
\label{sec:method_inject}

设第 $c$ 个局部冲突区的纵向范围为 $[s_c^-,s_c^+]$，预测占用管与规划路径重叠的时间区间为 $I_{cr}=[t_{cr}^{in},t_{cr}^{out}]$。将每个占用区间向两侧扩张时间安全量 $\eta_c$，并在规划时域 $[0,T]$ 内合并，则可安全到达该冲突区的时间窗集合为
\begin{equation}
 \mathcal W_c=[0,T]\setminus
 \bigcup_r[t_{cr}^{in}-\eta_c,t_{cr}^{out}+\eta_c].
 \label{eq:safe_windows}
\end{equation}
式\eqref{eq:safe_windows}同时产生在障碍物进入前通过和在其离开后通过两类时间同伦，与自车最早可达时间相交后得到可达窗 $W_{cq}=[\alpha_{cq},\beta_{cq}]$，其中 $\tau(s_c)$ 只用于候选排序，安全窗边界始终来自预测占用的补集。

对多个冲突点，MIKU 按 $s_c$ 递增建立分层有向图，节点 $(c,q)$ 表示在窗 $W_{cq}$ 内通过第 $c$ 个冲突点。若相邻选择的目标到达时间 $\hat t_{c-1}$、$\hat t_c$ 满足
\begin{equation}
 \hat t_c\geq \hat t_{c-1}+\frac{s_c-s_{c-1}}{v_{max}},
 \qquad \hat t_c\in W_{cq},
 \label{eq:temporal_edge}
\end{equation}
则在两节点间连边。节点代价度量 $|\hat t_c-\tau(s_c)|$，边代价处罚先通过/后通过类别的频繁切换，滚动规划时再对偏离上周期同伦的节点加入持续性代价。对分层图执行有界宽度的动态规划，得到同时满足多个冲突点的最优时间同伦序列。

选定窗 $W_{cq}=[\alpha_{cq},\beta_{cq}]$ 后，其语义被直接写入速度 QP 的位置上下边界。
\begin{equation}
 \begin{aligned}
 s_j^{ub}&=\min\!\left(s_j^{ub,0},\ s_c^- -\epsilon_s\right),
 && t_j<\alpha_{cq},\\
 s_j^{lb}&=\max\!\left(s_j^{lb,0},\ s_c^+ +\epsilon_s\right),
 && t_j\geq\beta_{cq}.
 \end{aligned}
 \label{eq:window_to_st}
\end{equation}
当窗与时域左端或右端相接时，式\eqref{eq:window_to_st}中对应的一侧约束为空，对所有冲突点的边界取交集即得最终 $\mathcal B_s$，若不存在完整时间同伦，则在第一个不可达冲突点前生成停车上界。该传递仍产生标准线性盒约束，因而速度 QP 的目标矩阵、运动学等式、凸性和稀疏结构保持不变。

\subsection{整体流程与复杂度}
\label{sec:method_algorithm}

算法\ref{alg:miku}给出从自车状态与障碍物预测轨迹到路径中心边界及速度时序约束的完整流程。

\begin{algorithm}[htbp]
\small
\caption{MIKU 多障碍物约束统一流程}
\label{alg:miku}
\textbf{输入}\quad $n$ 个障碍物及预测轨迹，道路物理边界，自车状态\par
\textbf{输出}\quad 中心路径边界 $\{[l^{-}(s),l^{+}(s)]\}$，ST 位置上下边界 $\mathcal B_s$\par
按式\eqref{eq:center_road}内缩道路物理边界，初始化所有截面\;
\For{每个障碍物 $i$}{
  估计其代表截面到达时间 $\tau(s_i^-)$，查询预测占据\;
  按式\eqref{eq:dynamic_delta}计算 $d_{buf,i}$，按式\eqref{eq:center_forbidden}得 $[u_i,v_i]$\;
}
按 $s_i^-$ 升序扫描，形成纵向连通分量\;
\For{每个连通分量 $C$}{
  删除与 $[L^-,L^+]$ 完全不相交的禁行区间\;
  组内按 $(u_i,-v_i)$ 排序，递推前缀最大值 $V_p$\;
  按式\eqref{eq:gap_def}计算 $g_0,\ldots,g_k$，保留宽度最大的 $K$ 个正宽候选带\;
}
按式\eqref{eq:spatial_homotopy}选择纵向连续的空间同伦序列\;
\For{每个选定的分量候选带}{
  固定组内绕行侧，对覆盖的每个 $s$ 查询 $\tau(s)$ 处的活跃禁行区间\;
  以式\eqref{eq:center_band}逐截面写回 $[l^-(s),l^+(s)]$\;
}
沿规划路径构造式\eqref{eq:prediction_tube}的障碍物预测占用管\;
按式\eqref{eq:safe_windows}生成各冲突点的安全时间窗\;
以式\eqref{eq:temporal_edge}连接分层时间同伦图并选择最优因果序列\;
按式\eqref{eq:window_to_st}将各选定窗与基础 ST 边界取交集\;
若启用交替到达时间更新，则由所得 $s(t)$ 重构并阻尼更新 $\tau(s)$，在达到迭代上限前且未达到收敛阈值时重新执行约束构造\;
\KwRet 中心路径边界与 ST 位置上下边界 $\mathcal B_s$
\end{algorithm}

设路径离散点数为 $K$、纵向连通分量数为 $m$、每层空间候选数为 $Q$，则扫描线分组与组内排序的总开销为 $O(n\log n)$，单截面候选间隙生成为 $O(k\log k)$，空间同伦动态规划为 $O(mQ^2)$。若时间同伦图每层保留至多 $B$ 个部分序列，平均可达窗数为 $R$ 且冲突点数为 $h$，其搜索开销为 $O(hBR)$，本文采用的 $Q=3$ 与 $B=8$ 均为有界在线开销。极端情况下，时变占据查询和朴素交互密度统计分别为 $O(nK)$ 与 $O(n^2)$，因此 MIKU 将指数级的单组方向组合转化为精确排序扫描，并以有界的空间与时间同伦搜索完成跨分量协调。

% ====== _parts/experiment.tex ======
\section{实验与分析}
\label{sec:experiment}

实验先以单截面最优宽度一致性、配对随机场景和消融比较总体性能，再以确定性场景和 Apollo 规划链接入定位具体机制，其中时间盲基线与 MIKU 共享场景、车辆约束及两阶段 QP，配对结果用于比较不同约束构造与传递方式的影响。

\subsection{仿真平台与场景设置}
\label{subsec:exp_setup}

实验在 Apollo 11.0 中进行，运行平台采用 Intel i9-14900HX 处理器、64\,GB 内存和 Arch Linux。输入包括平台场景、道路地图、障碍物消息与车辆模型，MIKU 仅替换 SL 路径边界和 ST 位置边界的构造过程，并沿用两阶段 QP 的目标函数、运动学等式与求解器。差异化额外安全裕度取 $[0.10,0.40]$\,m，交互密度半径取 10\,m，纵向分组则由障碍物占据区间是否重叠确定。

表\ref{tab:exp_scenarios}给出 8 个确定性场景，其中 P1--P4 为用于显露时间盲投影或局部选向失效的机制场景，分别覆盖横穿行人、动静混合障碍、窄路双侧约束和维修借道，C1--C4 则通过延长时间域或调整障碍物横向位置形成一一对应的常规可通行条件，用于比较新约束接入前后的无碰撞终点到达结果。

\begin{table}[htbp]
\centering
\caption{确定性机制场景的构型设置}
\label{tab:exp_scenarios}
\small
\begin{tabular}{c l p{7.4cm}}
\toprule
\textbf{场景对} & \textbf{构型} & \textbf{可比化调整与比较目的} \\
\midrule
P1 / C1 & 单行人横穿 & C1延长规划时域，观察动态占用下的先通过/后通过时序选择 \\
P2 / C2 & 行人加停车 & C2外移静态车辆横向中心，观察动静混合边界的分组与通行带构造 \\
P3 / C3 & 窄路加双侧交通锥 & C3将对称锥列改为轻微非对称使基线可通过，保留差异化裕度释放通行空间 \\
P4 / C4 & 单车道维修封闭加借道 & C4将水马横向中心调至贪心切换临界侧，观察组级选向的连续性 \\
\bottomrule
\end{tabular}
\end{table}

比较方法包括固定配置的 Apollo 解耦基线 B0、加入时变投影但保留局部贪心的 B1、执行 3 轮到达时间更新的 B2，以及固定离散精度的五维联合网格参考 B3，其中前三种方法与 MIKU 共享两阶段 QP，B3 用于对照通行能力与计算开销。

若轨迹在全时域内无碰撞且终点满足 $s_{\mathrm{end}}\geq s_{max}-1$\,m，则记为一次无碰撞终点到达。连续指标包括进度率、加加速度均方根（jerk RMS）和单次完整规划耗时，对同一随机种子生成的配对场景先计算方法间差值，再通过 5000 次百分位 bootstrap 给出 95\% 置信区间。无碰撞终点到达和碰撞事件等二值指标采用 McNemar 精确检验比较配对比例。

\subsection{最大间隙扫描的精确性与总体规划性能}
\label{subsec:exp_exactness}

本节首先在固定截面实例上比较第\ref{sec:method_band}节的 $O(k\log k)$ 最大间隙扫描与 $2^k$ 组合枚举最优解，检验二者能否得到相同的单截面最优宽度，随后通过配对随机场景比较各方法的安全性、通行能力、平顺性与在线开销，以评估 MIKU 的总体规划性能。对于任一固定横向截面，最优中心通行带宽度记为 $W=l^+-l^-$，其中车辆半宽和安全裕度已经计入道路可行边界与障碍物禁行边界，因而 $W$ 直接表示自车中心可放置区间的长度。

设 $W_{\mathrm{scan}}^{(r)}$ 和 $W_{\mathrm{enum}}^{(r)}$ 分别为第 $r$ 个实例中 $O(k\log k)$ 扫描与 $2^k$ 组合枚举得到的最优中心通行带宽度，两种实现的最大宽度差定义为
\begin{equation}
    \Delta W_{\max}=\max_r\left|W_{\mathrm{scan}}^{(r)}-W_{\mathrm{enum}}^{(r)}\right|.
    \label{eq:exp_width_difference}
\end{equation}
20 个固定随机种子各生成 200 组含 0--8 个禁行区间的固定截面实例，区间允许重叠、嵌套或越出道路边界，在所得 4,000 个实例中，$O(k\log k)$ 最大间隙扫描与 $2^k$ 组合枚举最优解得到的单截面中心通行带宽度逐例一致且 $\Delta W_{\max}=0$，表明该检验核对的是两种方法在固定截面区间模型下的最优宽度一致性，因而不能用于衡量整条轨迹的横向跟踪误差。

\label{subsec:exp_randomized}

完成单截面最优宽度一致性检验后，总体性能实验覆盖行人横穿、车辆切入、停车加对向来车、窄路多障碍物、交错动态障碍物、预测噪声和延迟横穿 7 个场景族，每族使用 500 个固定随机种子生成车辆与障碍物参数，共得到 3,500 个配对场景。所有方法接收同一预测场景，并以独立保留的真值场景检测碰撞，其中无碰撞终点到达率同时反映安全到达与任务完成，碰撞率衡量安全性，进度率描述未到达车辆的实际推进程度，jerk RMS 和完整规划耗时则分别反映轨迹平顺性与在线计算开销。

\begin{table}[htbp]
\centering
\caption{3,500 个配对随机场景的总体结果}
\label{tab:randomized_overall}
\small
\begin{tabular}{lrrrr}
\toprule
\textbf{指标} & \textbf{B0} & \textbf{B1} & \textbf{B2} & \textbf{MIKU} \\
\midrule
无碰撞终点到达率 (\%) & 62.7 & 52.4 & 62.1 & \textbf{77.5} \\
碰撞率 (\%) & 2.2 & 2.1 & 2.2 & \textbf{1.1} \\
平均进度率 (\%) & 77.5 & 71.1 & 76.6 & \textbf{86.4} \\
jerk RMS (m/s$^3$) & 1.252 & 1.315 & 1.300 & \textbf{1.132} \\
规划耗时 P50 (ms) & 10.11 & 11.37 & 31.20 & 9.98 \\
规划耗时 P95 (ms) & 56.79 & 63.24 & 169.49 & 55.93 \\
\bottomrule
\end{tabular}
\end{table}

表\ref{tab:randomized_overall}显示，MIKU 的无碰撞终点到达率较 B0 提高 14.8 个百分点，95\% 置信区间为 $[13.5,16.0]$ 且 McNemar 检验满足 $p<10^{-132}$，同时碰撞率下降 1.0 个百分点并使 jerk RMS 降低 0.120\,m/s$^3$。MIKU 与 B0 的规划耗时中位数均约为 10\,ms，而需要执行三轮规划的 B2 达到 31.20\,ms，说明安全到达率的提高没有增加典型调用开销。

作为同时离散搜索时间、纵向位置、横向位置和速度状态的高计算量联合搜索参考，B3 用于比较通行能力与计算开销。在每个场景族分层抽取 10 个种子后，MIKU 与 B3 在 70 个场景中的无碰撞终点到达率分别为 77.1\% 和 80.0\%，规划耗时中位数分别为 10.44\,ms 和 1,835.43\,ms，因而 MIKU 以约 $1/176$ 的中位耗时获得了接近固定离散网格参考的通行能力。

\subsection{典型场景中的失效与修复机制}
\label{subsec:exp_passrate}

配对随机场景用于估计各方法的总体性能差异，P1--P4 则通过路径上下界的变化定位基线失去可行性的具体机制，一一对应的 C1--C4 用于比较 MIKU 与基线在基线可通行构型中的无碰撞终点到达结果。表\ref{tab:randomized_overall}承担总体性能比较，P1--P4 与 C1--C4 提供边界层面的机制证据。

\begin{table}[htbp]
\centering
\caption{机制场景与对应对照场景的无碰撞终点到达结果}
\label{tab:exp_pass}
\small
\begin{tabular}{l p{6.2cm} cc}
\toprule
\textbf{场景组} & \textbf{比较目的} & \shortstack{\textbf{基线无碰撞}\\\textbf{终点到达数}} & \shortstack{\textbf{MIKU 无碰撞}\\\textbf{终点到达数}} \\
\midrule
机制场景 P1--P4 & 检查时间盲投影与局部选向造成的阻塞能否被修复 & 1/4 & 4/4 \\
对照场景 C1--C4 & 检查新约束是否保持原本可通行场景的规划能力 & 4/4 & 4/4 \\
\bottomrule
\end{tabular}
\end{table}

表\ref{tab:exp_pass}显示，基线与 MIKU 均能处理主要依靠速度时序避让的 P1 单行人横穿场景，但在引入静态障碍、双侧约束或密集施工障碍的 P2--P4 中，基线逐个确定绕行侧并依次更新边界，使部分障碍物抬高下界而另一些障碍物压低上界，最终因 $l^-(s)>l^+(s)$ 形成并非真实几何封闭造成的阻塞。MIKU 将纵向相关障碍物作为一组联合选择绕行侧，并从连续分割产生的候选间隙中构造纵向连续通行带，因此在 P1--P4 中均无碰撞到达终点。

在仅调整时间域或障碍物横向位置的 C1--C4 中，两种方法均实现 4 个场景的无碰撞终点到达，表明在这组配对构型下 MIKU 保持了基线的通行结果。图~\ref{fig:q34_narrow_sl}以 P3 窄路双侧约束为例给出边界层面的证据，其中基线在 $s\approx16$\,m 处因上下界交叉而失去可行域，MIKU 则通过组级选向保持连续通行带并完成规划。

\begin{figure}[htbp]
\centering
\begin{subfigure}[b]{0.48\textwidth}
  \centering
  \adjustbox{max width=\linewidth}{\input{../图片/fig_exp_03_narrow_cones_sl_a}}
  \caption{基线路径在 $s\approx16$\,m 处阻塞}
\end{subfigure}\hfill
\begin{subfigure}[b]{0.48\textwidth}
  \centering
  \adjustbox{max width=\linewidth}{\input{../图片/fig_exp_03_narrow_cones_sl_b}}
  \caption{MIKU 的连续通行带与路径}
\end{subfigure}
\caption{窄路多障碍物场景的 SL 路径边界与通行对比}
\caption*{P3 中基线的上下界交叉造成虚假阻塞，而 MIKU 的组级选向保留了连续通行带。}
\label{fig:q34_narrow_sl}
\end{figure}

\subsection{组件消融与灵敏度}
\label{subsec:exp_ablation}

表\ref{tab:exp_ablation_score}给出主实验同一组 3,500 个配对场景上的减法消融，其中 A1--A7 每次从完整 MIKU 中去除一项机制。

\begin{table}[htbp]
\centering
\caption{3,500 个配对场景上的减法消融结果}
\label{tab:exp_ablation_score}
\small
\begin{tabular}{clrrr}
\toprule
\textbf{变体} & \textbf{删除的机制} & \textbf{无碰撞终点到达率(\%)} & \textbf{碰撞率(\%)} & \textbf{进度率(\%)} \\
\midrule
A1 & 到达时间驱动投影 & 77.3 & 1.2 & 86.7 \\
A2 & 纵向连通分组 & 76.9 & 1.1 & 85.9 \\
A3 & 前 $K$ 个候选带的空间同伦 & 67.3 & 1.1 & 79.5 \\
A4 & 威胁度差异化裕度 & 66.7 & 1.1 & 78.9 \\
A5 & 时间同伦图 & 74.2 & 1.2 & 84.3 \\
A6 & 不确定性预测占用管 & 77.0 & 2.2 & 86.6 \\
A7 & 交替到达时间更新 & 76.6 & 1.0 & 85.9 \\
MIKU & 完整方法 & \textbf{77.5} & \textbf{1.1} & \textbf{86.4} \\
\bottomrule
\end{tabular}
\end{table}

去除 Top-K 空间同伦或差异化裕度后，无碰撞终点到达率分别下降 10.2 和 10.7 个百分点，表明在本配对场景中，密集通道的规划结果对多候选选向和按威胁度分配的边界较为敏感。去除时间同伦图使延迟横穿场景族的无碰撞终点到达率下降 22.8 个百分点，而去除不确定性预测占用管使预测噪声场景族的碰撞率上升 7.4 个百分点，这些差异分别对应时序选择和预测误差扩张的影响。

\label{subsec:exp_sensitivity}

在 P1--P4 上对 5 个权重施加 $\pm20\%$ 均匀扰动并重新归一化后，100 组固定样本的额外裕度绝对偏差均小于 0.019\,m，排序翻转仅出现在裕度差小于 0.01\,m 的近等价交通锥之间，说明权重扰动造成的边界变化仍处于厘米量级。

\subsection{Apollo 规划链中的典型行为}
\label{subsec:apollo_hil}

为考察同一约束构造在既有工程规划器中的接入表现，本文将 MIKU 写入 Apollo 11.0 的路径边界与 ST 位置边界，并在系统提供的场景、障碍物消息和车辆模型下观察超车、窄路通行与封闭导流三类行为，如图~\ref{fig:q34_apollo_hil}所示。

\begin{figure}[htbp]
\centering
\begin{subfigure}[b]{0.48\textwidth}
  \centering
  \adjustbox{max width=\linewidth}{\input{../图片/fig_eng_scn01_dreamview_b}}
  \caption{动态超车}
\end{subfigure}\hfill
\begin{subfigure}[b]{0.48\textwidth}
  \centering
  \adjustbox{max width=\linewidth}{\input{../图片/fig_eng_scn02_dreamview_b}}
  \caption{窄路通行}
\end{subfigure}\hfill
\par\medskip
\begin{subfigure}[b]{0.72\textwidth}
  \centering
  \adjustbox{max width=\linewidth}{\input{../图片/fig_eng_scn03_dreamview_b}}
  \caption{单车道导流}
\end{subfigure}
\caption{Apollo 规划链中的三类典型行为}
\caption*{依次对应借道超车、窄路居中和封闭导流通行。}
\label{fig:q34_apollo_hil}
\end{figure}

动态超车场景按 $\tau(s)$ 将低速车辆的预计占据写入横向边界，使路径阶段提前形成左侧借道带并在超越后切回原车道。窄路场景通过扫描线分组、最大间隙和差异化裕度选择中心通道，封闭导流场景则依靠分量内方向一致的空间同伦统一绕行侧，三类行为表明离散决策写入原 QP 边界后能够产生相应轨迹。

% ====== 5 结论 ======
\section{结论}
\label{sec:conclusion}

本文提出面向路径--速度解耦规划器的多障碍物时空同伦约束统一方法 MIKU。与时间盲投影和逐障碍物贪心选向相比，按到达时间构造时变通行带、并在组内作连续分割，扩大了路径层可交给速度规划的可行空间，且不更换原两阶段 QP。连续分割引理与最大间隙定理给出组内绕行的最优性条件，使 $2^k$ 组合搜索降为 $O(k\log k)$ 扫描。空间同伦的时序前提经安全窗和时间同伦图写成 ST 盒约束后，轨迹优化仍表述为分段加加速度路径 QP 与速度 QP。4,000 组随机区间、3,500 个配对场景以及 Apollo 规划链接入表明，该方法的无碰撞终点到达率高于时间盲基线，碰撞率更低，规划耗时中位数约为 10\,ms，并以约 $1/176$ 的中位耗时接近联合网格参考。后续工作将在感知--预测--控制闭环中检验同一约束传递。

\printbibliography[heading=bibintoc,title={参考文献}]

\end{document}
